\documentclass[a4paper,11pt]{article}

% ---------- Préambule minimal ----------
\usepackage[utf8]{inputenc}
\usepackage[T1]{fontenc}
\usepackage[french]{babel}
\usepackage{amsmath,amssymb}
\usepackage{enumitem}
\usepackage{tikz}
\usepackage{tasks}
\usetikzlibrary{arrows.meta}
\usepackage[margin=2.5cm]{geometry}

% Environnement "exercice" : \begin{exercice}{<numéro>} ... \end{exercice}
\newenvironment{exercice}[1]{%
  \par\bigskip\noindent\textbf{Exercice #1}\par\smallskip\noindent\ignorespaces
}{%
  \par
}

\begin{document}
{
  \normalsize
  \textbf{Seconde} \hfill \textbf{Vecteurs 1} \hfill \textbf{Fiche exos 2}
}

% =====================================================
\begin{exercice}{1}
Simplifier les expressions suivantes en utilisant la relation de Chasles.
\begin{tasks}(3)
\task $\overrightarrow{AD}+\overrightarrow{DE}$
\task $\overrightarrow{BC}+\overrightarrow{AB}$
\task $\overrightarrow{AB}-\overrightarrow{AC}$
\task $\overrightarrow{AD}+\overrightarrow{DE}+\overrightarrow{EI}$
\task $\overrightarrow{AB}-\overrightarrow{AC}-\overrightarrow{CB}$
\task $\overrightarrow{AB}-\overrightarrow{AC}+\overrightarrow{BA}-\overrightarrow{AC}$
\end{tasks}
\end{exercice}

% =====================================================
\begin{exercice}{2}
Représenter à chaque fois le vecteur demandé.

\begin{center}
\begin{tikzpicture}[scale=0.8]
  \def\w{4} \def\h{3}
  % grille fine sur toute la zone (3 colonnes x 3 lignes)
  \draw[gray!30,very thin,step=0.5] (0,0) grid (3*\w,3*\h);
  % séparations épaisses des 9 cases
  \foreach \x in {0,1,2,3} \draw[thick] (\x*\w,0) -- (\x*\w,3*\h);
  \foreach \y in {0,1,2,3} \draw[thick] (0,\y*\h) -- (3*\w,\y*\h);
  % ligne du haut (y entre 2h et 3h) : vecteur donné, 2u, 1/2 u
  \draw[-{Latex[length=2.5mm]},thick] (0.5*\w,1+2*\h) -- (0.75*\w,2.5*\h) node[midway,above left] {$\vec u$};
  \node at (\w+0.6,3*\h-0.5) {$2\vec u$};
  \node at (2*\w+0.6,3*\h-0.5) {$\tfrac12\vec u$};
  % ligne du milieu (y entre h et 2h) : -u, -3/2 u, 1/3 u
  \node at (0.6,2*\h-0.5) {$-\vec u$};
  \node at (\w+0.6,2*\h-0.5) {$-\tfrac32\vec u$};
  \node at (2*\w+0.6,2*\h-0.5) {$\tfrac13\vec u$};
  % ligne du bas (y entre 0 et h) : -2/3 u, 5/3 u, 1/6 u
  \node at (0.6,\h-0.5) {$-\tfrac23\vec u$};
  \node at (\w+0.6,\h-0.5) {$\tfrac53\vec u$};
  \node at (2*\w+0.6,\h-0.5) {$\tfrac16\vec u$};
\end{tikzpicture}
\end{center}
\end{exercice}

% =====================================================
\begin{exercice}{3}
On considère deux vecteurs $\vec u$ et $\vec v$. $A$ est un point donné.
\begin{enumerate}[label=\arabic*.]
\item Construire les points $B$ et $C$ tels que :
  \begin{tasks}[style=itemize](2)
  \task $\overrightarrow{AB}=\vec u+\vec v$
  \task $\overrightarrow{AC}=\vec u-\vec v$.
  \end{tasks}
\item Construire les points $P$, $Q$ et $R$ tels que :
  \begin{tasks}[style=itemize](3)
  \task $\overrightarrow{BP}=\dfrac23\vec u$
  \task $\overrightarrow{PQ}=-2\vec v$
  \task $\overrightarrow{QR}=\dfrac{-2}{3}\vec u$.
  \end{tasks}
\item Que constate-t-on ? Le justifier par un calcul vectoriel.
\end{enumerate}

\begin{center}
\begin{tikzpicture}[scale=0.8]
  \draw[gray!30,very thin,step=0.5] (0,0) grid (14,9);
  \draw[-{Latex[length=2.5mm]},thick] (4,5) -- (1,1) node[above left] {$\vec u$};
  \draw[-{Latex[length=2.5mm]},thick] (6,7) -- (9,7) node[above,midway] {$\vec v$};
  \node at (12,8) {$\times$};
  \node[right] at (12.15,8) {$A$};
\end{tikzpicture}
\end{center}
\end{exercice}

% =====================================================
\begin{exercice}{4}
$ABCD$ est un parallélogramme.
\begin{enumerate}[label=\arabic*.]
\item Placer les points $E$ et $F$ tels que :
  \begin{itemize}
  \item $\overrightarrow{AE}=2\overrightarrow{AB}$
  \item $\overrightarrow{BF}=\dfrac32\overrightarrow{AD}-\dfrac12\overrightarrow{AB}$.
  \end{itemize}
\item Montrer que $\overrightarrow{CE}=\overrightarrow{AB}-\overrightarrow{AD}$.
\item Montrer que $\overrightarrow{CF}=\dfrac{-1}{2}\overrightarrow{AB}+\dfrac12\overrightarrow{AD}$.
\item En déduire le réel $\lambda$ tel que $\overrightarrow{CE}=\lambda\overrightarrow{CF}$.
\end{enumerate}

\begin{center}
\begin{tikzpicture}
  \draw[gray!30,very thin,step=0.5] (0,0) grid (14,9);
  \coordinate (A) at (5+1,3+3);
  \coordinate (B) at (5+5,3+3);
  \coordinate (C) at (5+4,3+0);
  \coordinate (D) at (5+0,3+0);
  \draw (A) -- (B) -- (C) -- (D) -- cycle;
  \foreach \P/\lab/\pos in {A/A/above left, B/B/above right, C/C/below right, D/D/below left}
    \fill (\P) circle (1.5pt) node[\pos] {$\lab$};
\end{tikzpicture}
\end{center}
\end{exercice}

\newpage
% =====================================================
\begin{exercice}{5}
$ABC$ est un triangle.
\begin{enumerate}[label=\arabic*.]
\item Placer les points $M$ et $N$ définis par $\overrightarrow{BM}=2\overrightarrow{AC}$ et $\overrightarrow{AN}=\dfrac12\overrightarrow{BC}$.
\item Démontrer que $\overrightarrow{MC}=2\overrightarrow{CN}$.
\item Soient $I$ et $J$ les milieux respectifs des segments $[AB]$ et $[AC]$. Montrer que $\overrightarrow{BC}=2\overrightarrow{IJ}$.
\item Montrer que $ANJI$ est un parallélogramme.
\end{enumerate}

\begin{center}
\begin{tikzpicture}
  \draw[gray!30,very thin,step=0.5] (0,0) grid (16,9);
  \coordinate (A) at (5,2);
  \coordinate (C) at (7,6);
  \coordinate (B) at (13,1);
  \draw[thick] (A) -- (B) -- (C) -- cycle;
  \fill (A) circle (1.5pt) node[left] {$A$};
  \fill (B) circle (1.5pt) node[right] {$B$};
  \fill (C) circle (1.5pt) node[above] {$C$};
\end{tikzpicture}
\end{center}
\end{exercice}

% =====================================================
\begin{exercice}{6}
Sur la figure ci-dessous, $ABCD$ est un parallélogramme.

\begin{center}
\begin{tikzpicture}
  \draw[gray!30,very thin,step=0.5] (0,0) grid (12,9);
  \coordinate (A) at (2,4);
  \coordinate (B) at (7,4);
  \coordinate (C) at (9,8);
  \coordinate (D) at (4,8);
  \foreach \P/\lab/\pos in {A/A/left, B/B/below right, C/C/above right, D/D/above left}
    \fill (\P) circle (1.5pt) node[\pos] {$\lab$};
\end{tikzpicture}
\end{center}

\begin{enumerate}[label=\arabic*.]
\item Donner deux égalités de vecteurs de votre choix données par cette figure.
\item Placer les points $E$ et $F$ définis par :
  \begin{enumerate}[label=(\roman*)]
  \item $\overrightarrow{CF}=2\overrightarrow{CB}$
  \item $\overrightarrow{AE}=2\overrightarrow{AB}-\overrightarrow{BC}$
  \end{enumerate}
\item
  \begin{enumerate}[label=\alph*.]
  \item Exprimer $\overrightarrow{CF}$ en fonction de $\overrightarrow{CB}$ et $\overrightarrow{BF}$ puis montrer que $\overrightarrow{CB}=\overrightarrow{BF}$ en utilisant l'égalité (i).
  \item Que peut-on en déduire ?
  \end{enumerate}
\item
  \begin{enumerate}[label=\alph*.]
  \item À partir de l'égalité (ii), montrer que $\overrightarrow{AE}=\overrightarrow{AB}+\overrightarrow{AF}$ puis que $\overrightarrow{AB}=\overrightarrow{FE}$.
  \item Que peut-on en déduire ?
  \end{enumerate}
\item Montrer que le quadrilatère $CDFE$ est un parallélogramme.
\end{enumerate}
\end{exercice}

% =====================================================
\begin{exercice}{7}
Soit $ABC$ un triangle non aplati ($A$, $B$ et $C$ non alignés) et les points $D$ et $E$ tels que :
\[
\overrightarrow{AD}=\dfrac52\overrightarrow{AC}+\dfrac32\overrightarrow{CB}
\qquad\text{et}\qquad
\overrightarrow{CE}=-2\overrightarrow{AC}+\dfrac12\overrightarrow{AB}.
\]
\begin{enumerate}[label=\arabic*.]
\item Faire un dessin. Conjecturer le lien entre les points $B$, $D$ et $E$.
\item Nous allons démontrer la conjecture du point précédent.
  \begin{enumerate}[label=\alph*.]
  \item Exprimer $\overrightarrow{ED}$ en fonction des vecteurs $\overrightarrow{DA}$, $\overrightarrow{AC}$ et $\overrightarrow{CE}$, puis en fonction des seuls vecteurs $\overrightarrow{AB}$ et $\overrightarrow{AC}$.
  \item Exprimer $\overrightarrow{BD}$ en fonction des seuls vecteurs $\overrightarrow{AB}$ et $\overrightarrow{AC}$.
  \item On admet la propriété suivante :
  \begin{quote}
  \og Trois points $M$, $N$ et $P$ sont alignés ssi il existe un réel $\lambda$ tel que $\overrightarrow{MN}=\lambda\overrightarrow{MP}$. \fg{}
  \end{quote}
  Conclure.
  \end{enumerate}
\end{enumerate}
\end{exercice}

\end{document}